% Ch10 WS1 -- IMFs and the liquid state. Blocks 1-2.
\documentclass{apchem-worksheet}

\apcunitword{Chapter}
\apcunit{10}
\apcunittitle{Liquids and Solids}
\apcdoctitle{Worksheet 1 \textbullet\ Forces and Liquids}
\apcsubtitle{Zumdahl \S10.1--10.2 \textbullet\ compare size before comparing force type}

\begin{document}
\wsheader[Name the force AND the structural feature producing it. ``Stronger
IMFs'' restates the question --- say which force and why it is present.]

\problem[]{List every intermolecular force present:}
\begin{parts}
  \item \ce{CH3CH3}: \answerblank[6cm]{dispersion only}
  \item \ce{CH3Cl}: \answerblank[6cm]{dispersion $+$ dipole--dipole}
  \item \ce{CH3OH}: \answerblank[7.4cm]{dispersion, dipole--dipole, H-bonding}
  \item \ce{HF}: \answerblank[7.4cm]{dispersion, dipole--dipole, H-bonding}
\end{parts}

\problem[]{\ce{C2H5OH} boils at \SI{78}{\celsius} while \ce{CH3OCH3} boils
at \SI{-24}{\celsius}. Both have the formula \ce{C2H6O}. Explain the
\SI{100}{\celsius} gap. \answerlines[3]{Ethanol has an O--H group, so its
molecules hydrogen bond to one another. Dimethyl ether has no hydrogen
attached to oxygen, so only dipole--dipole and dispersion forces operate.
Hydrogen bonding is far stronger, so far more energy is needed to separate
ethanol molecules.}}

\problem[]{Zumdahl reports \ce{BI3} (dispersion only) boiling at
\SI{209.5}{\celsius} and \ce{NH3} (hydrogen bonding) at
\SI{-33.0}{\celsius}.}
\begin{parts}
  \item This appears to contradict ``hydrogen bonding is the strongest
        IMF.'' Resolve it. \answerlines[3]{Force \emph{type} does not set
        the ranking --- magnitude does. \ce{BI3} is a very large molecule
        with an enormous, highly polarizable electron cloud, so its
        dispersion forces are larger in absolute terms than the hydrogen
        bonds of tiny \ce{NH3}.}
  \item State the working rule this illustrates.
        \answerlines[2]{Compare molecular size and polarizability first when
        molecules differ greatly in size; compare force type only among
        molecules of comparable size.}
\end{parts}

\problem[]{Rank each set by increasing boiling point with a one-line reason:}
\begin{parts}
  \item \ce{Ne}, \ce{Ar}, \ce{Kr}
        \answerlines[2]{\ce{Ne} $<$ \ce{Ar} $<$ \ce{Kr} --- all have only
        dispersion, which grows with electron count.}
  \item \ce{CH4}, \ce{CH3F}, \ce{CH3OH}
        \answerlines[2]{\ce{CH4} $<$ \ce{CH3F} $<$ \ce{CH3OH} --- similar
        sizes, so force type decides: dispersion only, then dipole--dipole,
        then hydrogen bonding.}
\end{parts}

\problem[]{Surface tension, viscosity, capillary action:}
\begin{parts}
  \item Explain surface tension at the particle level.
        \answerlines[2]{Interior molecules are pulled equally in all
        directions, but surface molecules have no neighbours above and are
        pulled inward --- so the liquid minimizes its surface area.}
  \item Water forms a concave meniscus in glass; mercury forms a convex
        one. State which force dominates in each case.
        \answerlines[2]{Water: adhesive forces to glass exceed cohesive
        forces. Mercury: cohesive forces exceed adhesive forces to glass.}
  \item Give two reasons glycerol (\ce{C3H8O3}) is more viscous than
        ethanol. \answerlines[2]{It can form three hydrogen bonds per
        molecule rather than one, and its larger branched structure allows
        molecules to tangle, both resisting flow.}
\end{parts}

\problem[]{Which liquid in each pair has the higher vapour pressure at
\SI{25}{\celsius}? Justify.}
\begin{parts}
  \item Water or acetone (bp \SI{56}{\celsius}):
        \answerblank[5cm]{acetone --- weaker IMFs}
  \item Hexane or octane:
        \answerblank[9.0cm]{hexane --- fewer electrons, weaker dispersion}
\end{parts}

\begin{spiralstrip}{Chapters 8--9 \textbullet\ spiral}
  \item Hybridization of C in \ce{C2H4}: \answerblank[2cm]{sp$^2$}
  \item $\Delta H$ from bond energies:
        \answerblank[4cm]{broken $-$ formed}
  \item Higher lattice energy: \ce{KCl} or \ce{CaO}?
        \answerblank[2.4cm]{\ce{CaO}}
  \item $\sigma$ and $\pi$ in \ce{N2}:
        \answerblank[3.4cm]{1 $\sigma$, 2 $\pi$}
  \item Larger: \ce{S^2-} or \ce{Cl-}? \answerblank[2.2cm]{\ce{S^2-}}
\end{spiralstrip}

\end{document}
